\documentclass[../../../../SoftwareRequirementsSpecification.tex]{subfiles}
\begin{document}
% Richiede le macro definite in src/macros/useCaseMacros.tex
% (incluse nel preambolo di SoftwareRequirementsSpecification.tex)

\ucstart
\begin{tabularx}{\textwidth}{|l|X|}
  \hline
  \ucid{PT-05} \\ \hline
  \ucname{Assegnazione Scheda ad un Atleta} \\ \hline
  \ucactor{PT} \\ \hline
  \ucdescription{Il PT assegna ad un atleta una scheda
    già presente nel sistema.} \\ \hline
  \ucprecond{Il PT deve essere loggato e deve
    essere presente almeno una scheda configurata.} \\ \hline
  \ucpostcond{La scheda scelta risulta assegnata
    all'atleta, con i carichi della prima settimana specifici per
    lui.} \\ \hline
  \ucmainflow{
    \item Il PT si trova sulla pagina di dettaglio di un atleta,
      raggiunta tramite il caso d'uso PT-08.
    \item Il PT preme il tasto "Assegna nuova Scheda".
    \item Il sistema mostra la lista delle schede presenti nel sistema.
    \item Il PT preme sulla scheda che vuole usare.
    \item Il sistema mostra il form di inserimento dei dati
      aggiuntivi necessari, in particolare la data di inizio della
      scheda, e la durata della scheda, derivata dal numero di
      blocchi mensili configurati in PT-04; una volta inserita la
      data di inizio, il sistema mostra anche la data di fine che ne
      deriva (vedi le Note).
    \item Il PT inserisce i dati e preme il tasto "Continua".
    \item Il sistema mostra il form dei carichi degli esercizi delle
      sessioni della prima settimana. I campi richiesti dipendono
      dalla prescrizione di carico configurata in PT-04 (vedi le
      Note).
    \item Il PT inserisce i dati nel form e preme il tasto
      "Assegna Scheda".
    \item Il sistema assegna la Scheda all'Atleta.
    \item Il sistema invia una mail di notifica all'Atleta per
      informarlo che c'è una nuova scheda assegnata.
  } \\ \hline
  \ucaltflow{
    \textbf{4a. Creazione scheda ad hoc} \newline
    - Invece di scegliere una scheda già presente, il PT può crearne
    una da zero per l'atleta. \newline
    - Il flusso continua dal passo 1 del caso d'uso PT-03,
    successivamente dal passo 1 del caso d'uso PT-04. \newline
    - Dopo la creazione della scheda il flusso riprende dal
    passo 5. \newline
    \newline
    \textbf{6a. Sovrapposizione con una scheda già assegnata} \newline
    - Il sistema rileva che l'intervallo di date della nuova
    assegnazione si sovrappone a quello di una scheda già assegnata
    all'atleta. \newline
    - Il sistema mostra un messaggio di errore riportando il nome e
    l'intervallo di date della scheda in conflitto. \newline
    - Il flusso riprende dal passo 5. \newline
  } \\ \hline
  \ucnote{
    \textbf{Nota 1:} i carichi si fissano solo per le sessioni della
    prima settimana di allenamento; quelli delle settimane successive
    sono calcolati automaticamente dalle percentuali delle
    progressioni se queste sono presenti, altrimenti viene usato lo
    stesso carico per tutte le settimane. Cosa viene chiesto al PT
    dipende dalla prescrizione di carico dell'esercizio, fissata in
    PT-04:
    \begin{itemize}
      \item \textbf{carico assoluto}: il PT inserisce il peso;
      \item \textbf{percentuale del massimale}: il sistema mostra il
        massimale (1RM) dell'atleta per quell'esercizio, con la data
        in cui è stato rilevato, e il PT lo conferma o lo
        sovrascrive; il peso della prima settimana è calcolato
        applicandovi la percentuale. Se per quell'esercizio non è
        registrato alcun massimale il campo è vuoto ed
        obbligatorio;
      \item \textbf{corpo libero}: non viene chiesto nulla e per
        l'esercizio non esiste alcun carico.
    \end{itemize}
    Il massimale inserito qui è un dato di allenamento raccolto dal
    PT, non un dato anagrafico dell'atleta: resta associato
    all'atleta e viene sovrascritto ad ogni nuovo rilevamento. Il
    sistema non registra come il massimale sia stato
    misurato.\newline
    \textbf{Nota 2:} nel caso di esercizi con variazione drop-set
    viene inserito un solo carico: il peso delle ripetizioni in drop
    è ottenuto applicandovi il moltiplicatore configurato nella
    variante.\newline
    \textbf{Nota 3:} l'intervallo di date occupato da
    un'assegnazione va dalla data di inizio all'ultimo giorno
    dell'ultimo blocco mensile: ciascun blocco dura 4 settimane
    (28 giorni), quindi la data di fine è determinata dalla data di
    inizio e dal numero di blocchi, e non viene chiesta al PT. Due
    assegnazioni contigue, in cui una termina il giorno prima che
    l'altra inizi, non si sovrappongono e sono quindi ammesse.\newline
    \textbf{Nota 4:} i carichi fissati qui non cambiano se il
    massimale dell'atleta viene aggiornato in seguito: restano quelli
    calcolati al momento dell'assegnazione. Per adeguare una scheda
    già in corso il PT usa il caso d'uso PT-09 o assegna una nuova
    scheda.
  } \\ \hline
\end{tabularx}
\end{document}
